\documentclass[11pt]{article}
\usepackage{preamble}
\setlength{\parskip}{4pt}

\begin{document}

\daybanner{AP Statistics \textperiodcentered\ Unit 4 \textperiodcentered\ Topic 4.7 \textperiodcentered\ B: 2027-01-25 \textperiodcentered\ E: 2027-03-12 \textperiodcentered\ TEACHER KEY}{Two Samples of Service Times}

\sectionbanner{CED TETHER}

{\small
\begin{itemize}[leftmargin=1.5em,itemsep=2pt,topsep=2pt]
  \item \textbf{Skill 1.D} --- Identify an appropriate inference method for confidence intervals.
  \item \textbf{Skill 4.C} --- Verify that inference procedures apply in a given situation.
  \item \textbf{Skill 3.D} --- Construct a confidence interval, provided conditions for inference are met.
  \item \textbf{EU UNC-4} --- An interval of values should be used to estimate parameters, in order to account for uncertainty.
  \item \textbf{LO UNC-4.V} --- Identify an appropriate confidence interval procedure for a difference of two population means. [Skill 1.D]
  \item \textbf{EK UNC-4.V.1} --- Consider a simple random sample from population 1 of size $n_1$, mean $\mu_1$, and standard deviation $\sigma_1$ and a second simple random sample from population 2 of size $n_2$, mean $\mu_2$, and standard deviation $\sigma_2$. If the distributions of populations 1 and 2 are normal or if both $n_1$ and $n_2$ are greater than 30, then the sampling distribution of the difference of means, $\bar{x}_1$ - $\bar{x}_2$, is also normal. The mean for the sampling distribution of $\bar{x}_1$ - $\bar{x}_2$ is $\mu_1$ - $\mu_2$. The standard deviation of $\bar{x}_1$ - $\bar{x}_2$ is $\surd$($\sigma_1^2$/$n_1$ + $\sigma_2^2$/$n_2$).
  \item \textbf{EK UNC-4.V.2} --- The appropriate confidence interval procedure for one quantitative variable for two independent samples is a two-sample t-interval for a difference between population means.
  \item \textbf{LO UNC-4.W} --- Verify the conditions to calculate confidence intervals for the difference of two population means. [Skill 4.C]
  \item \textbf{EK UNC-4.W.1} --- In order to calculate confidence intervals to estimate a difference of population means, we must check for independence and that the sampling distribution is approximately normal:
  \item (a) To check for independence: (i) Data should be collected using two independent, random samples or a randomized experiment. (ii) When sampling without replacement, check that $n_1$ $\leq$ 10\% $N_1$ and $n_2$ $\leq$ 10\% $N_2$.
  \item (b) To check that the sampling distribution of ($\bar{x}_1$ - $\bar{x}_2$) should be approximately normal (shape): (i) If the observed distributions are skewed, both $n_1$ and $n_2$ should be greater than 30.
  \item \textbf{LO UNC-4.X} --- Determine the margin of error for the difference of two population means. [Skill 3.D]
  \item \textbf{EK UNC-4.X.1} --- For the difference of two sample means, the margin of error is the critical value (t*) times the standard error (SE) of the difference of two means.
  \item \textbf{EK UNC-4.X.2} --- The standard error for the difference in two sample means with sample standard deviations $s_1$ and $s_2$ is $\surd$($s_1^2$/$n_1$ + $s_2^2$/$n_2$).
  \item \textbf{LO UNC-4.Y} --- Calculate an appropriate confidence interval for a difference of two population means. [Skill 3.D]
  \item \textbf{EK UNC-4.Y.1} --- The point estimate for the difference of two population means is the difference in sample means, $\bar{x}_1$ - $\bar{x}_2$.
  \item \textbf{EK UNC-4.Y.2} --- For a difference of two population means where the population standard deviations are not known, the confidence interval is ($\bar{x}_1$ - $\bar{x}_2$) $\pm$ t* $\surd$($s_1^2$/$n_1$ + $s_2^2$/$n_2$) where $\pm$t* are the critical values for the central C\% of a t-distribution with appropriate degrees of freedom that can be found using technology.
\end{itemize}
}
\textbf{Essential question:} How does the sampling design determine whether a paired or independent-means confidence interval is trustworthy?\par
\textbf{DOK-3 skill:} 4.C \textperiodcentered\ \textbf{FRQ pattern:} {\small\texttt{design-\allowbreak{}selects-\allowbreak{}two-\allowbreak{}mean-\allowbreak{}interval}}\par
\textbf{DOK rationale:} Part (c) links the original sampling design to competing standard errors and intervals, diagnoses response-based pairing, and traces how false precision changes the population conclusion.\par

\frameworkphaseheader{Do Now (first take) $\rightarrow$ Explore (video + follow-along) $\rightarrow$ Exit (finish + turn in)}{1 $\rightarrow$ 3}{45}{%
\begin{itemize}[leftmargin=1.5em,itemsep=2pt,topsep=2pt]
  \item Display the board and ask for one sentence using the study description.
  \item Begin the 7.6 video follow-along at minute 5.
  \item Return to the board at minute 33. Students finish the shortened parts and justify their judgment.
  \item Collect at the door. Score part (c) only using the three required elements.
\end{itemize}
}{%
\begin{itemize}[leftmargin=1.5em,itemsep=2pt,topsep=2pt]
  \item Write a one-sentence first take before the video.
  \item Complete the video follow-along worksheet.
  \item Finish (a)--(c), revise the first take in the exit line, and turn in the sheet.
\end{itemize}
}{%
\begin{itemize}[leftmargin=1.5em,itemsep=2pt,topsep=2pt]
  \item Which specific number or study detail supports your choice?
  \item What claim would go beyond the evidence, and why?
\end{itemize}
}{Circulate and ask for evidence. Accept a concise response; do not require omitted calculations or a second full inference write-up.}

\begin{callout}{\IconWarn\ WATCH FOR}{calloutred}
\begin{itemize}[leftmargin=1.5em,itemsep=2pt,topsep=2pt]
  \item Choosing a speaker without a reason.
  \item Copying numbers without explaining what they support.
\end{itemize}
\end{callout}

\sectionbanner{SCORING PART (c) --- E / P / I}

\textbf{Expected elements}
\begin{itemize}[leftmargin=1.5em,itemsep=2pt,topsep=2pt]
  \item \textbf{uses-design} (required) --- Selects independent-sample analysis from the separate SRSs
  \item \textbf{sorting} (required) --- Explains that sorting observed records does not create design-based pairing and understates uncertainty
  \item \textbf{compares-intervals} (required) --- Contrasts zero inside the valid interval with zero outside the artificial paired interval
\end{itemize}

{\small\renewcommand{\arraystretch}{1.25}
\noindent\begin{tabular}{|p{0.5in}|p{5.9in}|}
\hline
\textbf{E} & Justifies the independent procedure, explains the artificial precision, and connects both intervals to the mean-difference conclusion. \\ \hline
\textbf{P} & Shows substantive valid reasoning for at least one required element, but does not meet all E requirements. Credit correct reasoning even when another claim is wrong. \\ \hline
\textbf{I} & Shows no substantive valid reasoning for any required element; a name, unsupported choice, or copied number alone does not count. \\ \hline
\end{tabular}}

\clearpage
\focusbanner{TODAY'S PROBLEM --- annotated}

Suppose two depots independently take SRSs of 36 weekday service-time records from monthly populations of 900 North and 1,200 South records. Technology gives $df=69.31$ and $t^*=1.995$ for a 95\% two-sample interval.\par\smallskip \textbf{Tariq:} ``I sorted both samples and paired equal ranks. The resulting differences have $\bar d=2.2$ and $s_d=3.0$, giving $(1.18,3.22)$. Pairing makes the estimate more precise.''\par \textbf{Mei:} ``The records were sampled separately, so I would use a two-sample interval. Sorting observed responses may change the uncertainty rather than reveal design-based pairs.''\par
\begin{center}
\textbf{Service-time samples (min)}\par\smallskip
\begin{tabular}{|l|c|c|c|c|}
\hline
\textbf{Depot} & \textbf{$N$} & \textbf{$n$} & \textbf{$\bar x$} & \textbf{$s$} \\ \hline
\textbf{North} & 900 & 36 & 42.3 & 8.4 \\ \hline
\textbf{South} & 1,200 & 36 & 40.1 & 7.6 \\ \hline
\end{tabular}
\end{center}
\begin{firsttakebox}{FIRST TAKE (5 min)}
Does sorting the times make the North and South records real pairs? Explain in one sentence.\par
\answer{Ungraded. Everyday language is enough before the video; formal vocabulary belongs in the finish.}
\end{firsttakebox}

\textit{Rules callout ``LET THE DESIGN CHOOSE THE INTERVAL (keep on the page)'' is printed on the student sheet.}\par\medskip

\dokbadge{1}~\textbf{(a)} Find the North-minus-South difference in sample mean service times.\par
\answer{$42.3-40.1=2.2$ minutes.}

\dokbadge{2}~\textbf{(b)} The independent-sample standard error is 1.88797 minutes. Use $t^*=1.995$ to calculate the 95\% interval.\par
\answer{$2.2\pm1.995(1.88797)=2.2\pm3.7665=(-1.5665,5.9665)$ minutes.}

\dokbadge{3}~\textbf{(c)} In 3--4 sentences, choose a procedure using the original sampling design. Compare the two intervals and explain how sorting the times could mislead a reader about whether the population means differ.\par
\answer{Mei is justified because the two samples were selected independently, without linked records. Sorting observed times does not create a paired study; the resulting $SE=3/\sqrt{36}=0.5$ makes the interval artificially narrow. The independent-sample interval $(-1.57,5.97)$ includes 0, so equal monthly population means remain plausible. The sorted-pair interval $(1.18,3.22)$ excludes 0 and could falsely suggest evidence that North has the higher population mean.}

\begin{summaryexitbox}{EXIT}
One thing the video changed about my first take: \hrulefill
\end{summaryexitbox}

\end{document}
